% EM structure and teaching revision 20260928
\subsection{组装能、相互作用能与自能}
把若干点电荷从无穷远逐个移入，先前已放置的电荷对后来电荷做功。每对电荷只计一次，因此
\[
 W=\sum_{i<j}\frac{kq_iq_j}{r_{ij}}
 =\frac12\sum_i q_i\varphi_{\mathrm{other}}(\boldsymbol r_i).
\]
右式的电势必须排除该点电荷自身的发散项。两个同号电荷由很远处靠近时 $W$ 增大，需要外力做正功；异号电荷靠近则相反。
连续分布可写为 $W=\tfrac12\int\rho\varphi\dd V$。对于半径 $R$ 的均匀实心带电球，也可逐层组装：已放好的内球电荷为 $q(r)=Qr^3/R^3$，新薄壳为 $\dd q=3Qr^2\dd r/R^3$，故
\[
 W=\int_0^R\frac{kq(r)}r\dd q=\frac{3kQ^2}{5R}.
\]
这包含整个连续球的自能；若题目仅问两部分之间的相互作用能，则应减去各部分各自的自能。把各区域的场直接平方相加会漏掉交叉项。
用于粒子运动时，先界定哪些带电体固定、哪些能运动。单粒子在固定外场中可用 $K+q\varphi=\mathrm{const}$；若场源也运动，则须把场源的动能和体系相互作用能一起计入。

